Put \(h_P(x):=\sum_{a\in\mathcal A}\omega_a(x)\mu_a(x)\).  Bounded weights, bounded means, the variance range, and \(e_{t,a}^{D,P}\geq\varepsilon\) imply that all quantities below are square integrable.  For any fixed \(t\), iterated conditional expectation and the non-anticipating assignment rule give
\[
\begin{aligned}
E_{\mathbb P_n^{D,P}}\!\left[
 \sum_{a\in\mathcal A}
 \frac{\mathbf 1\{A_t=a\}}{e_{t,a}^{D,P}}
 \omega_a(X_t)Y_t
 \ \middle|\ X_t,\mathcal F_{t-1}^{D,P}
\right]
&=\sum_{a\in\mathcal A}\omega_a(X_t)
 E_P[Y_t(a)\mid X_t]\\
&=h_P(X_t).
\end{aligned}
\]
The first equality uses consistency \(Y_t=Y_t(A_t)\), conditional randomization from the induced-experiment definition, and positivity for the division.  Since \(X_t\sim P_X\) by latent i.i.d. sampling, a second expectation yields
\[
\theta(P)=E_{\mathbb P_n^{D,P}}\!\left[
 \sum_{a\in\mathcal A}
 \frac{\mathbf 1\{A_t=a\}}{e_{t,a}^{D,P}}
 \omega_a(X_t)Y_t
\right].
\]
Thus the right-hand side is a computable observed-law functional because \(D\), hence its realized predictable probabilities, is known.  This proves identification; it does not define the causal target by the observed functional.

For the efficiency calculation, fix a static measurable allocation \(e\) with \(e_a(x)\geq\varepsilon\) and \(\sum_a e_a(x)=1\).  Work in the ambient dominated observed-data model through \(P\), with the treatment kernel held fixed.  Its canonical gradient for the displayed linear functional is
\[
\phi_e(O_t;P)
=h_P(X_t)-\theta(P)
 +\sum_{a\in\mathcal A}
 \frac{\mathbf 1\{A_t=a\}}{e_a(X_t)}
 \omega_a(X_t)\{Y_t-\mu_a(X_t)\}.
\]
Indeed, \(E_P[\phi_e]=0\).  Against a covariate-law score, the first term represents the derivative of \(E_P h_P(X)\); against an arm-specific conditional-outcome score, conditional expectation of the corresponding inverse-probability term is the derivative of \(E_P\{\omega_a(X)\mu_a(X)\}\); and against an assignment score the derivative is zero because the target does not depend on the known treatment kernel.  These three score spaces are mutually orthogonal, so this gradient is canonical.  This is an ambient-model efficiency statement; at a boundary point of a fixed Hölder ball or variance interval it should not be read as asserting the existence of a two-sided submodel contained in the restricted class.

Conditional centering makes the first term orthogonal to every residual term, and residual terms from distinct arms have zero product.  Therefore
\[
\begin{aligned}
\operatorname{Var}_P\{\phi_e(O_t;P)\}
&=\operatorname{Var}_P\{h_P(X_t)\}
 +E_P\!\left[
  \sum_{a\in\mathcal A}
  \frac{\omega_a(X_t)^2\sigma_a^2(X_t)}{e_a(X_t)}
 \right]\\
&=\operatorname{Var}_P\{h_P(X_t)\}
 +E_P\!\left[\sum_{a\in\mathcal A}\frac{m_a(X_t)}{e_a(X_t)}\right]
=V(e;P).
\end{aligned}
\]
Finally, the positivity-constrained simplex is compact pointwise, the objective is continuous in \(e\), and a measurable minimizer can be selected explicitly by the finite-dimensional convex program
\[
\min\left\{\sum_{a\in\mathcal A}\frac{m_a(x)}{u_a}:
u_a\geq\varepsilon,\ \sum_{a\in\mathcal A}u_a=1\right\}
\]
with deterministic lexicographic tie breaking.  Integrating its value gives exactly the infimum named \(V^*(P)\) in the oracle-benchmark definition.  The case in which some or all \(m_a(x)\) vanish is covered by the same compact program; nonuniqueness affects neither its value nor the conclusion.